\documentclass[11pt]{elegantbook}

\title{Cosmology Notes}
\subtitle{Theory, Observations and Methods}

\author{Hongfei Yu}
\institute{Department of Astronomy, University of Science and Technology of China}
\date{\today}
\version{0.1}
\bioinfo{Academic Page}{\url{https://phiyu.github.io}}

% \extrainfo{\textcolor{red}{\bfseries Caution: This template will no longer be maintained since January 1st, 2023. However, due to its large user base, maintenance has resumed and the template has been re-released as of 2026.}}

\logo{uphi.png}
\cover{cover.jpeg}

% modify the color in the middle of titlepage
\definecolor{customcolor}{RGB}{129, 92, 148}
\colorlet{coverlinecolor}{customcolor}
\usepackage{cprotect}

\addbibresource[location=local]{reference.bib} % bib

\begin{document}

\maketitle

\frontmatter
\tableofcontents

\mainmatter

% 
% ========== Part I: Fundamental Cosmology ==========
\part{Fundamental Cosmology}

% Ch: Cosmological Principle and FLRW Metric
\chapter{Cosmological Principle and FLRW Metric}

\section{Cosmological Principle}

\section{FLRW Metric}

% Ch: Friedmann Equations and LCDM Model
\chapter{Friedmann Equations and LCDM Model}

\section{Friedmann Equations}
\subsection{First Friedmann Equation}

\subsection{Second Friedmann Equation}

\section{LCDM Model}

\subsection{Cosmological Constant and Dark Energy}

\subsection{Dark Matter}


% ========== Part II: Statistical Cosmology ==========
\part{Statistical Cosmology}

% Ch: Perturbation Theory
\chapter{Perturbation Theory}

\section{Linear Perturbation Theory}

\section{Feynman Diagram and Kernels}

\section{Lagrangian Perturbation Theory}



% Ch: Correlation Function and Spectrum
\chapter{Correlation Function and Spectrum}

\section{Correlation Function}

\section{Power Spectrum and Bispectrum}
\subsection{Power Spectrum}

\subsection{Bispectrum}

\subsection{Consistency Relation}


\subsection{Squeezed Limit}


% Ch: Halo Model and HOD
\chapter{Halo Model and HOD}

\section{Spherical Collapse Model}

\section{Press-Schechter Formalism}

\section{Progenitor Distribution}

\section{Halo Occupation Distribution}

% ========== Part III: Observational Cosmology ==========
\part{Observational Cosmology}

% Ch: Galaxy Clustering and RSD
\chapter{Galaxy Clustering and RSD}

\section{Galaxy Clustering}

\section{Bias and Stochasticity}

\section{Redshift Space Distortion}

% Ch: CMB, BAO and SN
\chapter{CMB, BAO and Supernova}

\section{Cosmic Microwave Background}
\subsection{Temperature Anisotropy}

\subsection{Polarization Anisotropy}

\subsection{Sache-Wolfe Effect}

\subsection{Sunyaev-Zel'dovich Effect}

\section{Baryon Acoustic Oscillation}


\section{Supernova}

% Ch: Lensing
\chapter{Lensing}

\section{Gravitational Lensing}

\subsection{Weak Lensing}

\subsection{Strong Lensing}

\section{CMB Lensing}

% Ch: Other Probes
\chapter{Other Probes}

\section{Cosmic Chronometer}

\section{$21\,\mathrm{cm}$ Cosmology}

\section{Potential Decay Rate}

% ========== Part IV: Computational Cosmology ==========
\part{Computational Cosmology}

% Ch: N-body Simulation
\chapter{$N$-body Simulation}

\section{Initial Conditions}

\section{Equations of Motion}


\section{Evolution Algorithms}


\section{Particle Assignment}

% Ch: Bayesian Statistics and Monte Carlo Methods
\chapter{Bayesian Statistics and Monte Carlo Methods}

\section{Bayesian Statistics}

\section{Monte Carlo Methods}

\section{Monte Carlo Markov Chain}

% Ch: Emulators and Machine Learning
\chapter{Emulators and Machine Learning}
\section{Gaussian Processes}

\section{Convolutional Neural Networks}

% ========== Part V: Advanced Topics in Cosmology ==========
\part{Advanced Topics in Cosmology}

% Ch: Early Universe
\chapter{Early Universe}
\section{Inflation}

\section{Primordial Black Holes}

\section{Primordial Gravitational Waves}

% Ch: Dark Energy and Modified Gravity
\chapter{Dark Energy and Modified Gravity}

\section{Dynamical Dark Energy}

\section{Clustering Dark Energy}

\section{Interacting Dark Energy}

\section{Modified Gravity Theories}

% Ch: Constraints on Parameters and Models
\chapter{Constraints on Parameters and Models}

\section{Hubble Constant and $S_8$ Tensions}

\section{Neutrino Mass and Hierarchy}

\section{Non-Gaussianity}

% Ch: Reconstruction of the Universe
\chapter{Reconstructing the Universe}

\section{Reconstruction of Density Field}

\section{Reconstruction of Velocity Field}

\section{Reconstruction of Initial Conditions}


\end{document}